\documentclass{cslab}

\usepackage{listings}
\usepackage{xcolor}

% Python code styling
\lstset{
    language=Python,
    basicstyle=\ttfamily\small,
    keywordstyle=\color{blue}\bfseries,
    stringstyle=\color{red},
    commentstyle=\color{green!50!black}\itshape,
    columns=fullflexible,
    showstringspaces=false,
    breakatwhitespace=false,
    breaklines=true,
    frame=single,
    tabsize=4
}

% Lab information
\labnumber{}
\labtitle{Capstone Lab: From Pixels to Pipelines}
\course{CSCI 128: Coding for Problem Solving}
\courseshort{CSCI 128}
\semester{Winter 2026}
\timelimit{2 hours}
\logoimage{LogoWordsBottom.png}

\begin{document}

\maketitle

\begin{objectives}
    \item Use variables, types, and expressions to compute values
    \item Load images and manipulate pixels using the RGB colour model
    \item Apply colour filters (grayscale, negative, brightness, tint)
    \item Perform geometric transformations (mirror, crop, scale)
    \item Compose a multi-image collage onto a canvas
    \item Read and filter a CSV dataset
    \item Connect image processing with data analysis in a single pipeline
\end{objectives}

% -----------------------------------------------------------
\section{Setup}

\begin{setup}
\textbf{Before you begin:}
\begin{enumerate}
    \item Create a folder: \texttt{csci128/capstone}
    \item Inside it, create a subfolder called \texttt{data} and another called \texttt{output}
    \item Create \texttt{data/cities.csv} with the following content exactly:
\end{enumerate}

\begin{lstlisting}[language={},numbers=none,frame=single]
name,country,population,area_km2,coastal
Halifax,Canada,403000,5490,yes
Sydney,Australia,5312000,12368,yes
Lyon,France,522000,518,no
Nairobi,Kenya,4397000,696,no
Vancouver,Canada,2463000,2878,yes
Zurich,Switzerland,415000,87,no
Lima,Peru,10719000,2672,yes
Oslo,Norway,693000,454,yes
Kampala,Uganda,3652000,189,no
Quebec City,Canada,800000,3154,no
\end{lstlisting}

Save the file and confirm it appears at \texttt{capstone/data/cities.csv} before continuing.
\end{setup}

% -----------------------------------------------------------
\section{Lab Tasks}

% -----------------------------------------------------------
\task{Warming Up with Variables and Types}

Before writing any image or CSV code, let's make sure the concepts from the very first
lectures are solid.

\textbf{Create \texttt{warmup.py}:}

\begin{lstlisting}
# -------------------------------------------------------
# Part A: variables and arithmetic
# -------------------------------------------------------

# A standard HD image is 1920 x 1080 pixels.
# Store the width and height in variables.
width  = ___
height = ___

# Calculate and print the total number of pixels.
total_pixels = ___
print("Total pixels in an HD image:", ___)

# Each pixel stores three colour channels (R, G, B).
# Each channel is one byte (0-255).
# How many bytes does one HD image occupy?
bytes_per_image = ___
print("Bytes per image:", ___)
print("Kilobytes per image:", ___)

# -------------------------------------------------------
# Part B: types and conversion
# -------------------------------------------------------

# The string below came from reading a CSV file.
raw_population = "403000"

# Convert it to an integer and store it.
population = ___

# Calculate and print: how many people per square km if
# the city covers 5490 km^2?
area = 5490
density = ___
print(f"Population density: {density:.1f} people/km2")

# -------------------------------------------------------
# Part C: string formatting
# -------------------------------------------------------

city_name = "Halifax"
country   = "Canada"

# Print one line that reads:
#   Halifax, Canada --- population: 403,000
print(___)
\end{lstlisting}

\begin{note}
Fill in every \texttt{\_\_\_} blank with the correct expression or value.  There are no
imports needed --- pure Python arithmetic and strings.
\end{note}

% -----------------------------------------------------------
\task{Colour Filters on a City Photo}

For this task you need \textbf{one image} of your choice --- a photo of a city, a
landscape, or anything with colour.  Name it \texttt{data/photo.jpg}.

\textbf{Create \texttt{filters.py}:}

\begin{lstlisting}
from PIL import Image


def grayscale(img):
    """Return a grayscale copy of img.

    Use the standard luminance formula to convert each pixel to gray.
    """
    result = img.copy()
    pixels = result.load()
    w, h = result.size
    for x in range(w):
        for y in range(h):
            r, g, b = pixels[x, y]
            # --- YOUR CODE HERE ---
            pass
    return result


def negative(img):
    """Return a negative copy of img."""
    result = img.copy()
    pixels = result.load()
    w, h = result.size
    for x in range(w):
        for y in range(h):
            r, g, b = pixels[x, y]
            # --- YOUR CODE HERE ---
            pass
    return result


def adjust_brightness(img, amount):
    """Return a copy of img with brightness shifted by amount.

    Parameters
    ----------
    img    : PIL Image
    amount : int   positive = brighter, negative = darker
    """
    result = img.copy()
    pixels = result.load()
    w, h = result.size
    for x in range(w):
        for y in range(h):
            r, g, b = pixels[x, y]
            # --- YOUR CODE HERE ---
            pass
    return result


def tint(img, tr, tg, tb):
    """Return a copy of img with a colour tint applied.

    Each channel should be scaled by its corresponding factor
    (tr, tg, tb), then clamped to [0, 255].
    """
    result = img.copy()
    pixels = result.load()
    w, h = result.size
    for x in range(w):
        for y in range(h):
            r, g, b = pixels[x, y]
            # --- YOUR CODE HERE ---
            pass
    return result


# -------------------------------------------------------
# Apply and save all four filters
# -------------------------------------------------------
original = Image.open("data/photo.jpg")

grayscale(original).save("output/photo_gray.jpg")
negative(original).save("output/photo_negative.jpg")
adjust_brightness(original, 50).save("output/photo_bright.jpg")
adjust_brightness(original, -50).save("output/photo_dark.jpg")
tint(original, 1.0, 0.7, 0.7).save("output/photo_warm.jpg")

print("All filters saved to output/")
\end{lstlisting}

\begin{warning}
Each function must work on a \textbf{copy} of the image (\texttt{img.copy()}) so the
original is never modified.  This lets you chain filters without side effects.
\end{warning}

\begin{hint}
Clamping a value \texttt{v} to the range \texttt{[0, 255]} is a one-liner:
\begin{lstlisting}[numbers=none,frame=none]
v = max(0, min(255, v))
\end{lstlisting}
\end{hint}

% -----------------------------------------------------------
\task{Geometric Transformations}

\textbf{Create \texttt{transforms.py}:}

\begin{lstlisting}
from PIL import Image


def mirror_horizontal(img):
    """Return a left-right mirrored copy of img."""
    result = img.copy()
    src    = img.load()
    dst    = result.load()
    w, h   = img.size
    for x in range(w):
        for y in range(h):
            # --- YOUR CODE HERE ---
            pass
    return result


def crop(img, x1, y1, x2, y2):
    """Return the rectangular region [x1:x2, y1:y2] of img.

    Parameters
    ----------
    x1, y1 : int   top-left corner of the crop region
    x2, y2 : int   bottom-right corner (exclusive)
    """
    new_w = x2 - x1
    new_h = y2 - y1
    result = Image.new("RGB", (new_w, new_h))
    src = img.load()
    dst = result.load()
    for x in range(new_w):
        for y in range(new_h):
            # --- YOUR CODE HERE ---
            pass
    return result


def scale_down(img, factor):
    """Return a copy of img scaled down by an integer factor.

    Parameters
    ----------
    factor : int   e.g. 2 means half the size in each dimension
    """
    w, h   = img.size
    new_w  = w // factor
    new_h  = h // factor
    result = Image.new("RGB", (new_w, new_h))
    src    = img.load()
    dst    = result.load()
    for x in range(new_w):
        for y in range(new_h):
            # --- YOUR CODE HERE ---
            pass
    return result


# -------------------------------------------------------
# Apply transformations
# -------------------------------------------------------
original = Image.open("data/photo.jpg")

mirror_horizontal(original).save("output/photo_mirror.jpg")

# Crop the centre quarter of the image
w, h = original.size
crop(original, w // 4, h // 4, 3 * w // 4, 3 * h // 4).save(
    "output/photo_crop.jpg"
)

scale_down(original, 2).save("output/photo_half.jpg")

print("Transformations saved to output/")
\end{lstlisting}

\begin{note}
For \texttt{mirror\_horizontal} you are reading from \texttt{src} (the \emph{original}
pixel array) and writing to \texttt{dst} (the \emph{copy}).  If you read and write the
same pixel array you will corrupt the image halfway through the loop.
\end{note}

% -----------------------------------------------------------
\task{Reading and Filtering the Cities Dataset}

\textbf{Create \texttt{read\_cities.py}:}

\begin{lstlisting}
import csv


def load_cities(filepath):
    """Load cities.csv and return a list of dicts."""
    with open(filepath, newline="") as f:
        # --- YOUR CODE HERE ---
        # Use csv.DictReader to read the file.
        # Return a list of all rows.
        pass


def filter_coastal(cities):
    """Return only the cities where coastal == 'yes'."""
    # --- YOUR CODE HERE ---
    pass


def filter_by_min_population(cities, min_pop):
    """Return cities whose population is >= min_pop."""
    # --- YOUR CODE HERE ---
    pass


def sort_by_population(cities, descending=True):
    """Return cities sorted by population."""
    # --- YOUR CODE HERE ---
    pass


def print_city_table(cities):
    """Print a simple table of city name, country, and population."""
    print(f"{'City':<15} {'Country':<15} {'Population':>12}")
    print("-" * 44)
    for city in cities:
        print(f"{city['name']:<15} {city['country']:<15} "
              f"{int(city['population']):>12,}")


# -------------------------------------------------------
# Main
# -------------------------------------------------------
cities = load_cities("data/cities.csv")

print("=== All cities ===")
print_city_table(cities)

print("\n=== Coastal cities only ===")
print_city_table(filter_coastal(cities))

print("\n=== Cities with population >= 1,000,000 ===")
big_cities = filter_by_min_population(cities, 1_000_000)
print_city_table(sort_by_population(big_cities))
\end{lstlisting}

\begin{hint}
\texttt{sorted()} does not modify the original list --- it returns a new sorted list.
To sort by a specific column, you will need to pass a \textbf{key function} to
\texttt{sorted()}.  Remember that numeric columns read from a CSV are strings and
must be converted before comparing.
\end{hint}

% -----------------------------------------------------------
\task{Generating City Badges}

For each city in a filtered list you will create a small coloured image \textbf{badge}
that displays the city name and its population density.  This is where data analysis
meets image generation.

\textbf{Create \texttt{badges.py}:}

\begin{lstlisting}
import csv
from PIL import Image, ImageDraw, ImageFont


BADGE_W = 300
BADGE_H = 150

# Colour scheme: coastal cities get a blue badge, inland cities get green.
COASTAL_COLOUR = (70, 130, 180)    # steel blue
INLAND_COLOUR  = (60, 140,  80)    # forest green
TEXT_COLOUR    = (255, 255, 255)   # white


def compute_density(population, area_km2):
    """Return population density rounded to one decimal place."""
    # --- YOUR CODE HERE ---
    pass


def make_badge(city_name, country, density, is_coastal):
    """Return a BADGE_W x BADGE_H PIL Image for one city.

    The badge must show:
      - A filled background (blue for coastal, green for inland)
      - The city name in large text at the top
      - The country name in smaller text below the city name
      - The density (people/km2) at the bottom
    """
    bg_colour = COASTAL_COLOUR if is_coastal else INLAND_COLOUR
    img  = Image.new("RGB", (BADGE_W, BADGE_H), bg_colour)
    draw = ImageDraw.Draw(img)

    # --- YOUR CODE HERE ---
    # Draw the city name, country, and density onto the badge.

    return img


# -------------------------------------------------------
# Load data and generate one badge per city
# -------------------------------------------------------
with open("data/cities.csv", newline="") as f:
    cities = list(csv.DictReader(f))

badge_images = []   # collect badges in this list

for city in cities:
    name     = city["name"]
    country  = city["country"]
    pop      = int(city["population"])
    area     = float(city["area_km2"])
    coastal  = city["coastal"] == "yes"

    density  = compute_density(pop, area)
    badge    = make_badge(name, country, density, coastal)

    badge.save(f"output/badge_{name.replace(' ', '_')}.jpg")
    badge_images.append(badge)
    print(f"Badge created: {name}")

print(f"\n{len(badge_images)} badges saved to output/")
\end{lstlisting}

\begin{note}
\texttt{ImageDraw.Draw(img)} gives you a drawing surface attached to \texttt{img}.
Anything you draw on it modifies \texttt{img} in place, so there is no need to return
the draw object.  The \texttt{draw.text((x, y), text, fill=colour)} method places
\texttt{text} with its top-left corner at \texttt{(x, y)}.
\end{note}

% -----------------------------------------------------------
\task{Assembling the Collage}

Now bring everything together: load the badges produced in Task 5, arrange them in a
grid, and save the final collage.  Then write a summary CSV of the cities that appear
in the collage.

\textbf{Create \texttt{collage.py}:}

\begin{lstlisting}
import csv
import os
from PIL import Image


BADGE_W   = 300
BADGE_H   = 150
COLS      = 5          # number of columns in the grid
PADDING   = 10         # pixels of space between badges
BG_COLOUR = (30, 30, 30)   # dark background


def build_collage(badge_paths, cols, badge_w, badge_h, padding, bg):
    """Arrange badge images in a grid and return the canvas.

    Parameters
    ----------
    badge_paths : list of str    paths to individual badge images
    cols        : int            number of columns
    badge_w, badge_h : int       expected size of each badge
    padding     : int            gap between badges (pixels)
    bg          : tuple          background RGB colour

    Returns
    -------
    PIL Image
    """
    rows = (len(badge_paths) + cols - 1) // cols   # ceiling division

    canvas_w = cols * badge_w + (cols + 1) * padding
    canvas_h = rows * badge_h + (rows + 1) * padding
    canvas   = Image.new("RGB", (canvas_w, canvas_h), bg)

    for i, path in enumerate(badge_paths):
        badge = Image.open(path)

        col = i % cols
        row = i // cols

        # --- YOUR CODE HERE ---
        # Calculate the x and y position for this badge based on its
        # row and column, then paste it onto the canvas.
        pass

    return canvas


def save_summary_csv(badge_paths, output_path):
    """Write a CSV listing the cities that appear in the collage."""
    with open(output_path, "w", newline="") as f:
        writer = csv.DictWriter(f, fieldnames=["position", "city", "filename"])
        writer.writeheader()
        for i, path in enumerate(badge_paths):
            filename = os.path.basename(path)
            # Extract city name: "badge_Quebec_City.jpg" -> "Quebec City"
            city_name = filename.replace("badge_", "").replace(".jpg", "")
            city_name = city_name.replace("_", " ")
            writer.writerow({
                "position": i + 1,
                "city": city_name,
                "filename": filename,
            })


# -------------------------------------------------------
# Collect badge paths (sorted alphabetically for a tidy grid)
# -------------------------------------------------------
badge_paths = sorted([
    os.path.join("output", f)
    for f in os.listdir("output")
    if f.startswith("badge_") and f.endswith(".jpg")
])

if not badge_paths:
    print("No badges found in output/ - run badges.py first!")
else:
    collage = build_collage(badge_paths, COLS, BADGE_W, BADGE_H,
                            PADDING, BG_COLOUR)
    collage.save("output/collage.jpg")
    print(f"Collage saved: {collage.size[0]}x{collage.size[1]} pixels, "
          f"{len(badge_paths)} badges")

    save_summary_csv(badge_paths, "output/collage_summary.csv")
    print("Summary CSV saved to output/collage_summary.csv")
\end{lstlisting}

\begin{note}
Run the scripts \textbf{in order}:
\begin{enumerate}
    \item \texttt{warmup.py} --- no dependencies
    \item \texttt{filters.py} --- needs \texttt{data/photo.jpg}
    \item \texttt{transforms.py} --- needs \texttt{data/photo.jpg}
    \item \texttt{read\_cities.py} --- needs \texttt{data/cities.csv}
    \item \texttt{badges.py} --- needs \texttt{data/cities.csv}
    \item \texttt{collage.py} --- needs the badge files produced by \texttt{badges.py}
\end{enumerate}
\end{note}


\end{document}
